\section{The No-Zeno engine (mechanized algebra)}\label{sec:no-zeno-engine}

This section isolates the \emph{pure inequality logic} behind our ``No-Zeno'' decision engine.
Nothing here is specific to Navier--Stokes: the same argument applies to \emph{any} multiscale
process whose refinement events admit (i) an increment lower bound and (ii) a divergence check.
All PDE content enters later only through attempting to justify the hypotheses (WORK/CAP/DIV).

\paragraph{Mechanization boundary.}
All results in this section are mechanized in Lean in a discrete setting:
crossing times are sequences $t:\N\to\R$ and hypotheses are stated using finite sums over $\N$.
Continuous-time integral inequalities are used later only as motivational analogues and are not
formalized in this paper.

\subsection{Crossing times and the Zeno condition}

Let $t:\N\to\R$ be an increasing sequence of \emph{crossing times}, where $t(n)$ is interpreted as
the first time a multiscale ``frontier'' reaches level $n$ (e.g.\ a shell index threshold).
Define increments
\[
\Delta t_n := t(n+1)-t(n)\qquad (n\in\N).
\]

\begin{definition}[Zeno cascade]\label{def:zeno}
We say a \emph{Zeno cascade} occurs if the crossing times are bounded above:
\[
\exists\,T\in\R\ \ \text{s.t.}\ \ \forall n\in\N,\ t(n)\le T.
\]
Equivalently, $\sup_n t(n)<\infty$, i.e.\ the frontier reaches arbitrarily fine levels in finite time.
\end{definition}

\subsection{Increment lower bounds imply No-Zeno}

The basic mechanism is: if each increment $\Delta t_n$ is bounded below by a nonnegative
quantity whose partial sums are unbounded, then $t(n)$ cannot remain bounded.

\begin{definition}[Increment certificate]\label{def:phi}
Let $\Phi:\N\to\R$ be a nonnegative sequence. We say $\Phi$ is an \emph{increment certificate}
for $t$ if for all $n$,
\[
\Phi(n)\ \le\ t(n+1)-t(n).
\]
\end{definition}

\begin{theorem}[Unbounded partial sums imply no finite upper bound]\label{thm:unbounded-sums-no-zeno}
Assume $\Phi$ is an increment certificate for $t$ and that its partial sums are unbounded:
\[
\forall M\in\R,\ \exists n\in\N\ \text{s.t.}\ M \le \sum_{k=0}^{n-1}\Phi(k).
\]
Then $t$ has no finite upper bound, hence no Zeno cascade occurs.
\end{theorem}

\begin{proof}[Proof idea]
The increment inequality telescopes to
$t(n)\ge t(0)+\sum_{k=0}^{n-1}\Phi(k)$.
If the partial sums on the right are unbounded, then $t(n)$ cannot be bounded above.
\end{proof}

\subsection{WORK/CAP/DIV form}

For the applications we care about, the increment certificate is produced from a
\emph{work/capacity} inequality.

\begin{definition}[WORK/CAP/DIV]\label{def:work-cap-div}
Let $w,\mathrm{Cap}:\N\to\R$ be sequences with $\mathrm{Cap}(n)>0$.
\begin{itemize}
\item \textbf{WORK:} $w(n)$ is the ``work quantum'' required to advance from level $n$ to $n+1$.
\item \textbf{CAP:} $\mathrm{Cap}(n)$ is the ``capacity'' available at level $n$ (a bound relating work
to input energy and/or throughput on the crossing interval).
\item \textbf{WORK$\times$CAP inequality:} for all $n$,
\[
w(n)\ \le\ \mathrm{Cap}(n)\,\bigl(t(n+1)-t(n)\bigr).
\]
\item \textbf{DIV:} the ratios $w(n)/\mathrm{Cap}(n)$ have unbounded partial sums:
\[
\forall M\in\R,\ \exists n\in\N\ \text{s.t.}\ M \le \sum_{k=0}^{n-1}\frac{w(k)}{\mathrm{Cap}(k)}.
\]
\end{itemize}
\end{definition}

\begin{theorem}[WORK/CAP $\Rightarrow$ increment lower bound]\label{thm:workcap-increment}
Assume $\mathrm{Cap}(n)>0$ and
\[
w(n)\ \le\ \mathrm{Cap}(n)\,(t(n+1)-t(n)).
\]
Then $w(n)/\mathrm{Cap}(n)\ \le\ t(n+1)-t(n)$.
\end{theorem}

\begin{theorem}[WORK+CAP+DIV $\Rightarrow$ No-Zeno]\label{thm:workcap-no-zeno}
Assume the WORK$\times$CAP inequality and DIV from Definition~\ref{def:work-cap-div}.
Then $t$ has no finite upper bound, hence no Zeno cascade occurs.
\end{theorem}

\subsection{Dyadic wrapper (``doubling'' language)}

Sometimes one only checks crossings along a dyadic subsequence (``each doubling of scale'').
Define a dyadic-sampled ratio
\[
\Phi_{\mathrm{dyad}}(j) := \frac{w(2^j)}{\mathrm{Cap}(2^j)}.
\]
A dyadic version of Theorem~\ref{thm:workcap-no-zeno} follows by applying the same logic to
the sampled times $j\mapsto t(j)$ together with the dyadic increment hypotheses.

\begin{remark}[Mechanization note]
All results in this section are mechanized in Lean in this repository:
\texttt{formal/Fluids/Zeno.lean} (increment certificates and ``no finite upper bound'' from unbounded sums)
and \texttt{formal/Fluids/ZenoWorkCap.lean} (the WORK/CAP division step and its dyadic wrapper), using Lean 4 and the mathlib ecosystem \citep{lean4,demoura2015lean,mathlib2020}.
\end{remark}
